\begin{abstract}

Neural networks trained with standard objectives exhibit behaviors characteristic of probabilistic inference: soft clustering, prototype specialization, and Bayesian uncertainty tracking. These phenomena appear across architectures---in attention mechanisms, classification heads, and energy-based models---yet existing explanations rely on loose analogies to mixture models or post-hoc architectural interpretation. We provide a direct explanation. For any objective with log-sum-exp structure over distances or energies, the gradient with respect to each distance is exactly the negative posterior responsibility of the corresponding component: $\partial L / \partial d_j = -r_j$. The identity is algebraic, requiring only differentiability; it is a specialization of Fisher's identity from the classical EM literature, and its significance here is its address: standard neural objectives instantiate it without modification. The immediate consequence is that gradient descent on such objectives performs a generalized expectation-maximization implicitly, with responsibilities arising as gradients to be applied rather than auxiliary variables to be computed. No explicit inference algorithm is required because inference is embedded in the optimization. This result unifies three regimes of learning under a single mechanism: unsupervised mixture modeling, where responsibilities are fully latent; attention, where responsibilities are conditioned on queries; and cross-entropy classification, where supervision clamps responsibilities to targets. Our claims live at training time: the responsibility-weighted gradient dynamics recently documented in transformers follow necessarily from the objective's geometry. The in-context Bayesian computation that trained transformers perform at inference time is the endpoint of these dynamics, not their per-step content. For this class of objectives, learning and inference are the same process viewed at different levels.

\end{abstract}
